%% kapa: exam=2 round=33 year=2022 period=1 subject=appraisal_practice subject_ko=감정평가실무 seq=1 number=1 points=40
%% source: Q-Net 감정평가사 기출문제 (cst00309) artlSeq=5210838 file=raw/1008/5210838/2022년 제33회 감정평가사 2차시험 1교시 문제지(감정평가실무).pdf
%% conversion: kordoc 4.24.51 -> tools/gen2.py (auto); review=proofread-gpt-astra; flags=none
\begin{longquestion}{【문제 1】}
감정평가사 甲은 A군수로부터 「도로법」에 따른 도로에 편입되는 토지ㆍ지장물 등에 대한 협의보상평가를 의뢰받았다. 관련 법규 및 이론에 의거 제시된 자료를 활용하여 다음의 각 물음에 답하시오. (40점) \par
\medskip\noindent 물음 1) 대상토지에 대한 보상액을 산정하시오. (단, 시산가액에 대한 합리성 검토는 하지 않음) (10점) \par
\medskip\noindent 물음 2) 소유자 乙은 아래 자료와 같이 주거용 건축물을 신축하여 현재까지 거주하고 있다. 지장물 중 건축물에 대한 보상액을 산정하시오. (단, 건축물은 이전이 불가능하며, 이주대책 등은 고려하지 않음) (10점) \par
\medskip\noindent 물음 3) 지장물 중 수목에 대한 보상평가방법을 설명하고, 소나무(관상수)에 대한 보상액을 산정하시오. (6점) \par
\medskip\noindent 물음 4) 소유자 乙씨 부부(부부와 함께 동거하던 아들 1명은 현재 징집으로 인한 입영 중임)는 주거용 건축물이 공익사업에 편입됨으로 인하여 추가적으로 지급받을 수 있는 보상으로 이주대책, 주거이전비, 이사비 등이 있다. 하지만 사업시행자는 이주대책을 별도로 수립ㆍ실시하지 않는다고 한다. 공익사업의 시행으로 주거용 건축물을 제공함에 따라 생활의 근거를 상실하게 된 소유자 乙이 수령할 수 있는 보상액은 각각 얼마인가? (8점) \par
\medskip\noindent 물음 5) 물건조서의 꿀벌(양봉), 닭(산란계)에 대하여 축산업의 손실에 대한 보상평가를 하려고 한다. 영업손실의 보상대상인 영업의 일반적인 요건은 성립하는 것으로 가정한다. 꿀벌(양봉), 닭(산란계)에 대하여 < 자료 10 >을 토대로 축산업 손실보상의 대상여부를 판단하여 보상액을 산정하시오. (6점) \par
< 자료 1 > 평가의뢰 내역 등 \par
1. 사업의 종류: ○○ ～ ○○간 도로건설공사 \par
2. 도로구역결정고시일: 2021. 12. 31. \par
3. 사업시행자: A군수 \par
4. 보상평가의뢰일자: 2022. 6. 1. \par
5. 제시된 가격시점: 2022. 7. 31. \par
6. 현장조사완료일: 2022. 6. 30. \par
7. 토지조서 및 대상토지특성 \par
{\small\centering\renewcommand{\arraystretch}{1.8}
\resizebox{\ifdim\width>\linewidth\linewidth\else\width\fi}{!}{
\begin{tabular}{|c|c|c|c|c|c|c|c|c|c|c|}
\hline
\multirow{2}{*}{\makecell{기호}} & \multirow{2}{*}{\makecell{소재지}} & \multirow{2}{*}{\makecell{지번}} & \multirow{2}{*}{\makecell{지목}} & \multirow{2}{*}{\makecell{현실\\이용\\상황}} & \multirow{2}{*}{\makecell{전체\\면적\\(m\textsuperscript{2})}} & \multirow{2}{*}{\makecell{편입\\면적\\(m\textsuperscript{2})}} & \multirow{2}{*}{\makecell{용도지역\\및 지구}} & \multicolumn{3}{|c|}{\makecell{비고}} \\
\cline{9-11}
 &  &  &  &  &  &  &  & \makecell{도로\\교통} & \makecell{형상\\지세} & \makecell{기타} \\
\hline
\makecell{1} & \makecell{A군\\B면\\C리} & \makecell{106번지} & \makecell{대} & \makecell{단독\\주택} & \makecell{330} & \makecell{330} & \makecell{계획관리\\지역} & \makecell{세로\\(가)} & \makecell{사다리형\\완경사} & \makecell{후면\\주택\\지대} \\
\hline
\end{tabular}}\par}

※ 토지 소유자: 乙 \par
8. 물건조서 및 대상건축물특성 \par
{\small\centering\renewcommand{\arraystretch}{1.25}
\resizebox{\ifdim\width>\linewidth\linewidth\else\width\fi}{!}{
\begin{tabular}{|c|c|c|c|c|c|c|c|c|}
\hline
\multirow{2}{*}{\makecell{기호}} & \multirow{2}{*}{\makecell{소재지}} & \multirow{2}{*}{\makecell{지번}} & \multirow{2}{*}{\makecell{물건의 종류}} & \multirow{2}{*}{\makecell{구조 및\\규격}} & \multirow{2}{*}{\makecell{수량}} & \multicolumn{3}{|c|}{\makecell{비고}} \\
\cline{7-9}
 &  &  &  &  &  & \makecell{사용승인일} & \makecell{등급} & \makecell{내용\\년수} \\
\hline
\makecell{1} & \makecell{A군\\B면 C리} & \makecell{106번지} & \makecell{단독주택} & \makecell{벽돌조 슬래\\브지붕 1층} & \makecell{88m\textsuperscript{2}} & \makecell{2017. 7. 1.} & \makecell{상급} & \makecell{50년} \\
\hline
\makecell{2} & \makecell{A군\\B면 C리} & \makecell{106번지} & \makecell{부속창고} & \makecell{벽돌조 슬래\\브지붕 1층} & \makecell{8m\textsuperscript{2}} & \makecell{2017. 7. 1.} & \makecell{중급} & \makecell{45년} \\
\hline
\makecell{3} & \makecell{A군\\B면 C리} & \makecell{106번지} & \makecell{야외화장실} & \makecell{벽돌조 슬래\\브지붕 1층} & \makecell{3m\textsuperscript{2}} & \makecell{2017. 7. 1.} & \makecell{중급} & \makecell{45년} \\
\hline
\makecell{4} & \makecell{A군\\B면 C리} & \makecell{106번지} & \makecell{소나무\\(관상수)} & \makecell{H3.5 ×\\W1.5 × R15} & \makecell{1주} &  &  &  \\
\hline
\makecell{5} & \makecell{A군\\B면 C리} & \makecell{106번지} & \makecell{꿀벌(양봉)} &  & \makecell{30군} &  &  &  \\
\hline
\makecell{6} & \makecell{A군\\B면 C리} & \makecell{106번지} & \makecell{닭(산란계)} & \makecell{70일령 이상} & \makecell{20마리} &  &  &  \\
\hline
\end{tabular}}\par}

※ 물건 소유자: 乙 \par
< 자료 2 > 표준지공시지가 \par
{\small\centering\renewcommand{\arraystretch}{1.25}
\resizebox{\ifdim\width>\linewidth\linewidth\else\width\fi}{!}{
\begin{tabular}{|c|c|c|c|c|c|c|c|c|c|c|}
\hline
\makecell{기호} & \makecell{소재지} & \makecell{지목} & \makecell{면적\\(m\textsuperscript{2})} & \makecell{이용\\상황} & \makecell{용도지역} & \makecell{도로\\교통} & \makecell{형상\\지세} & \makecell{공시\\지가\\(원/m\textsuperscript{2})} & \makecell{비고} & \makecell{공시\\기준일} \\
\hline
\makecell{가} & \makecell{A군\\B면 C리\\50번지} & \makecell{대} & \makecell{210} & \makecell{주상용} & \makecell{계획관리\\지역} & \makecell{세로\\(가)} & \makecell{부정형\\평지} & \makecell{400,000} & \makecell{계획도로저촉\\10\%} & \makecell{2021.\\1. 1.} \\
\hline
\makecell{나} & \makecell{A군\\B면 C리\\50번지} & \makecell{대} & \makecell{210} & \makecell{주상용} & \makecell{계획관리\\지역} & \makecell{세로\\(가)} & \makecell{부정형\\평지} & \makecell{450,000} & \makecell{계획도로저촉\\10\%} & \makecell{2022.\\1. 1.} \\
\hline
\makecell{다} & \makecell{A군\\B면 C리\\65번지} & \makecell{대} & \makecell{300} & \makecell{단독주택} & \makecell{계획관리\\지역} & \makecell{소로\\한면} & \makecell{사다리형\\평지} & \makecell{300,000} & \makecell{계획도로저촉\\25\%} & \makecell{2021.\\1. 1.} \\
\hline
\makecell{라} & \makecell{A군\\B면 C리\\65번지} & \makecell{대} & \makecell{300} & \makecell{단독주택} & \makecell{계획관리\\지역} & \makecell{소로\\한면} & \makecell{사다리형\\평지} & \makecell{330,000} & \makecell{계획도로저촉\\25\%} & \makecell{2022.\\1. 1.} \\
\hline
\end{tabular}}\par}

주 1) 상기 표준지공시지가는 공히 대상토지와 인근지역에 소재함 \par
2) 표준지공시지가 기호 가), 나)는 주택상가혼용지대이고, 표준지공시지가 기호 다), 라)는 주택지대임 \par
< 자료 3 > 거래사례 \par
{\small\centering\renewcommand{\arraystretch}{1.25}
\resizebox{\ifdim\width>\linewidth\linewidth\else\width\fi}{!}{
\begin{tabular}{|c|c|c|c|c|c|c|c|c|c|}
\hline
\makecell{기호} & \makecell{소재지} & \makecell{지목} & \makecell{면적\\(m\textsuperscript{2})} & \makecell{이용상황} & \makecell{용도지역} & \makecell{도로\\교통} & \makecell{형상\\지세} & \makecell{거래가액\\(원/m\textsuperscript{2})} & \makecell{거래\\일자} \\
\hline
\makecell{A} & \makecell{A군\\B면 C리\\42번지} & \makecell{대} & \makecell{160} & \makecell{상업용} & \makecell{계획관리지역} & \makecell{소로\\한면} & \makecell{사다리형\\평지} & \makecell{700,000} & \makecell{2021.\\12. 1.} \\
\hline
\makecell{B} & \makecell{D군\\E면 F리\\15번지} & \makecell{대} & \makecell{270} & \makecell{단독주택} & \makecell{계획관리지역} & \makecell{세로\\(가)} & \makecell{사다리형\\완경사} & \makecell{390,000} & \makecell{2021.\\12. 1.} \\
\hline
\end{tabular}}\par}

주 1) 거래사례는 대상토지 및 표준지공시지가와 동일수급권 유사지역 또는 인근지역에 소재하며, 공히 지역요인은 대등함 \par
2) 거래사례는 거래당사자간의 사정이 개입되지 않은 정상적인 거래로 판단됨 \par
3) 거래사례 중 기호 A)는 주택상가혼용지대이고, 기호 B)는 후면 주택지대임 \par
< 자료 4 > 지가변동률 \par
{\small\centering\renewcommand{\arraystretch}{1.25}
\resizebox{\ifdim\width>\linewidth\linewidth\else\width\fi}{!}{
\begin{tabular}{|c|c|c|}
\hline
\makecell{기간} & \makecell{A군 계획관리지역\\지가변동률(\%)} & \makecell{비고} \\
\hline
\makecell{2021. 01. 01 ～\\2021. 12. 31(누계)} & \makecell{8.340} & \makecell{2021. 12. 01 ～\\2021. 12. 31: 0.150\%임} \\
\hline
\makecell{2022. 01. 01 ～\\2022. 05. 31(누계)} & \makecell{5.270} &  \\
\hline
\makecell{2022. 05. 01 ～\\2022. 05. 31} & \makecell{0.132} &  \\
\hline
\end{tabular}}\par}
\par{\small\centering\renewcommand{\arraystretch}{1.25}
\resizebox{\ifdim\width>\linewidth\linewidth\else\width\fi}{!}{
\begin{tabular}{|c|c|c|}
\hline
\makecell{기간} & \makecell{D군 계획관리지역\\지가변동률(\%)} & \makecell{비고} \\
\hline
\makecell{2021. 01. 01 ～\\2021. 12. 31(누계)} & \makecell{7.450} & \makecell{2021. 12. 01 ～\\2021. 12. 31: 0.130\%임} \\
\hline
\makecell{2022. 01. 01 ～\\2022. 05. 31(누계)} & \makecell{5.070} &  \\
\hline
\makecell{2022. 05. 01 ～\\2022. 05. 31} & \makecell{0.145} &  \\
\hline
\end{tabular}}\par}

주 1) 지가변동률은 용도지역별 지가변동률을 적용하며, 생산자물가지수는 고려하지 않기로 함 \par
2) 지가변동률은 2022년 6월 이후는 고시되지 않아서 5월 지가변동률을 연장ㆍ추정하여 적용함 \par
3) 지가변동률은 백분율로서 소수점 넷째자리에서 반올림하여 셋째자리까지 표시함 \par
< 자료 5 > 토지 가치형성요인 비교자료 \par
1. 접근조건 \par
1) 대상토지는 표준지 가), 나), 다), 라) 대비 각각 20\% 열세함 \par
2) 표준지 가), 나), 다), 라)는 거래사례 A) 와 대등, 거래사례 B) 대비 20\% 우세함 \par
2. 환경조건 \par
1) 대상토지는 표준지 가), 나) 대비 30\%, 표준지 다), 라) 대비 10\% 각각 열세함 \par
2) 표준지 가), 나)는 거래사례 A)와 대등, 거래사례 B) 대비 30\% 우세함 \par
3) 표준지 다), 라)는 거래사례 A) 대비 30\% 열세, 거래사례 B) 대비 10\% 우세함 \par
3. 기타 격차율 자료 \par
1) 토지이용상황 \par
{\small\centering\renewcommand{\arraystretch}{1.25}
\resizebox{\ifdim\width>\linewidth\linewidth\else\width\fi}{!}{
\begin{tabular}{|c|c|c|c|}
\hline
\makecell{구분} & \makecell{주거용} & \makecell{주상용} & \makecell{상업용} \\
\hline
\makecell{주거용} & \makecell{1.00} & \makecell{1.10} & \makecell{1.30} \\
\hline
\makecell{주상용} & \makecell{0.91} & \makecell{1.00} & \makecell{1.18} \\
\hline
\makecell{상업용} & \makecell{0.77} & \makecell{0.85} & \makecell{1.00} \\
\hline
\end{tabular}}\par}

2) 형상 \par
{\small\centering\renewcommand{\arraystretch}{1.25}
\resizebox{\ifdim\width>\linewidth\linewidth\else\width\fi}{!}{
\begin{tabular}{|c|c|c|}
\hline
\makecell{구분} & \makecell{사다리형} & \makecell{부정형} \\
\hline
\makecell{사다리형} & \makecell{1.00} & \makecell{0.95} \\
\hline
\makecell{부정형} & \makecell{1.05} & \makecell{1.00} \\
\hline
\end{tabular}}\par}
3) 경사 \par
{\small\centering\renewcommand{\arraystretch}{1.25}
\resizebox{\ifdim\width>\linewidth\linewidth\else\width\fi}{!}{
\begin{tabular}{|c|c|c|}
\hline
\makecell{구분} & \makecell{평지} & \makecell{완경사} \\
\hline
\makecell{평지} & \makecell{1.00} & \makecell{0.91} \\
\hline
\makecell{완경사} & \makecell{1.10} & \makecell{1.00} \\
\hline
\end{tabular}}\par}

4) 도로접면 \par
{\small\centering\renewcommand{\arraystretch}{1.25}
\resizebox{\ifdim\width>\linewidth\linewidth\else\width\fi}{!}{
\begin{tabular}{|c|c|c|}
\hline
\makecell{구분} & \makecell{세로가} & \makecell{소로한면} \\
\hline
\makecell{세로가} & \makecell{1.00} & \makecell{1.18} \\
\hline
\makecell{소로한면} & \makecell{0.85} & \makecell{1.00} \\
\hline
\end{tabular}}\par}

5) 도시ㆍ군계획시설 \par
{\small\centering\renewcommand{\arraystretch}{1.25}
\resizebox{\ifdim\width>\linewidth\linewidth\else\width\fi}{!}{
\begin{tabular}{|c|c|c|}
\hline
\makecell{구분} & \makecell{일반} & \makecell{도로} \\
\hline
\makecell{일반} & \makecell{1.00} & \makecell{0.85} \\
\hline
\makecell{도로} & \makecell{1.18} & \makecell{1.00} \\
\hline
\end{tabular}}\par}

※ 도시ㆍ군계획시설에 대한 요인비교치는 소수점이하 셋째자리에서 반올림하여 둘째자리까지 표시함 \par
< 자료 6 > 건축물 재조달원가 관련 자료 \par
1. 표준단가 \par
(단위: 원/m\textsuperscript{2}) \par
{\small\centering\renewcommand{\arraystretch}{1.25}
\resizebox{\ifdim\width>\linewidth\linewidth\else\width\fi}{!}{
\begin{tabular}{|c|c|c|c|c|c|}
\hline
\makecell{용도} & \makecell{구조} & \makecell{상급} & \makecell{중급} & \makecell{하급} & \makecell{내용년수} \\
\hline
\makecell{단독주택} & \makecell{벽돌조 슬래브 지붕} & \makecell{1,400,000} & \makecell{1,200,000} & \makecell{1,000,000} & \makecell{50년} \\
\hline
\makecell{창고} & \makecell{벽돌조 슬래브 지붕} & \makecell{400,000} & \makecell{360,000} & \makecell{320,000} & \makecell{45년} \\
\hline
\makecell{화장실} & \makecell{벽돌조 슬래브 지붕} & \makecell{1,300,000} & \makecell{1,100,000} & \makecell{900,000} & \makecell{45년} \\
\hline
\end{tabular}}\par}

※ 대상건축물 및 주거용 건축물 거래사례에도 동일하게 적용함 \par
(단위: 원/m\textsuperscript{2}) \par
2. 단독주택 부대설비 보정단가 \par
{\small\centering\renewcommand{\arraystretch}{1.25}
\resizebox{\ifdim\width>\linewidth\linewidth\else\width\fi}{!}{
\begin{tabular}{|c|c|c|}
\hline
\makecell{구분} & \makecell{위생ㆍ급배수설비} & \makecell{난방설비} \\
\hline
\makecell{보정단가} & \makecell{30,000} & \makecell{70,000} \\
\hline
\end{tabular}}\par}

※ 대상건축물 및 주거용 건축물 거래사례에도 동일하게 적용함 \par
< 자료 7 > 주거용 건축물 거래사례 자료 등 \par
1. 주거용 건축물 거래사례 \par
{\small\centering\renewcommand{\arraystretch}{1.5}
\resizebox{\ifdim\width>\linewidth\linewidth\else\width\fi}{!}{
\begin{tabular}{|c|c|c|c|c|c|c|c|}
\hline
\makecell{소재지} & \makecell{건축물\\면적\\(m\textsuperscript{2})} & \makecell{이용\\상황} & \makecell{구조} & \makecell{사용\\승인일} & \makecell{등급} & \makecell{거래가액\\거래일자} & \makecell{비고} \\
\hline
\multirow{3}{*}{\makecell{A군\\B면\\C리\\200번지}} & \makecell{88} & \makecell{단독\\주택} & \makecell{벽돌조\\슬래브지붕 1층} & \makecell{2017. 6. 1.} & \makecell{상급} & \multirow{3}{*}{\makecell{130,000,000원\\2022. 6. 1.}} & \multirow{3}{*}{\makecell{건축물\\일체의\\거래임}} \\
\cline{2-6}
 & \makecell{8} & \makecell{부속\\창고} & \makecell{벽돌조\\슬래브지붕 1층} & \makecell{2017. 6. 1.} & \makecell{중급} &  &  \\
\cline{2-6}
 & \makecell{3} & \makecell{야외\\화장실} & \makecell{벽돌조\\슬래브지붕 1층} & \makecell{2017. 6. 1.} & \makecell{중급} &  &  \\
\hline
\end{tabular}}\par}

주 1) 거래사례는 토지를 수반하지 않은 주거용 건축물만의 거래로서 보상대상 건축물과 유사하여 비교가능성이 높은 것으로 판단됨 \par
2) 거래사례는 인근지역 내의 거래당사자간의 사정이 개입되지 않은 정상적인 거래로 판단됨 \par
2. 건축물 가치형성요인 비교자료 \par
가. 대상건축물은 주거용 건축물 거래사례 대비 현상 및 관리상태에서 5\% 우세함 \par
나. 건축비는 연간 6\% 증가함 \par
※ 건축비 상승률을 고려한 지수는 소수점 셋째자리에서 절사하여 둘째자리까지 표시함 \par
다. 면적 등 기타 제반 가치형성요인은 대등함 \par
< 자료 8 > 수목 보상평가 참고자료 \par
1. 이식에 소요되는 비용은 다음과 같이 조사됨 \par
(단위: 주당) \par
{\small\centering\renewcommand{\arraystretch}{1.25}
\resizebox{\ifdim\width>\linewidth\linewidth\else\width\fi}{!}{
\begin{tabular}{|c|c|c|c|c|c|c|c|}
\hline
\makecell{구분} & \makecell{굴취비} & \makecell{운반비} & \makecell{상하차비} & \makecell{식재비} & \makecell{재료비} & \makecell{부대비용} & \makecell{소계} \\
\hline
\makecell{이식비} & \makecell{70,000} & \makecell{50,000} & \makecell{30,000} & \makecell{130,000} & \makecell{30,000} & \makecell{20,000} & \makecell{330,000} \\
\hline
\end{tabular}}\par}

※ 이식은 물리적으로 가능한 것으로 판단됨 \par
2. 수목의 취득가격은 400,000원/주 임 \par
3. 고손율을 적용할 경우 20\% 임 \par
4. 감수율을 적용할 경우 이식1차년: 100\%, 이식2차년: 80\%, 이식3차년: 40\% 임 \par
< 자료 9 > 이사비 등 보상 관련 참고자료 \par
1. 가구원수에 따른 1년분의 평균생계비는 1인당 15,106,000원 임 \par
※ 가구원수에 따른 1년분의 평균생계비 = 농가경제조사통계의 연간 전국평균 가계지출비 ÷ 가구당 전국평균 농가인구 \par
2. 도시근로자가구의 가구원수별 월평균 명목가계지출비 \par
{\small\centering\renewcommand{\arraystretch}{1.25}
\resizebox{\ifdim\width>\linewidth\linewidth\else\width\fi}{!}{
\begin{tabular}{|c|c|c|c|}
\hline
\makecell{구분} & \makecell{1인 가구} & \makecell{2인 가구} & \makecell{3인 가구} \\
\hline
\makecell{월평균 가계지출비} & \makecell{2,277,700원} & \makecell{3,334,200원} & \makecell{4,665,400원} \\
\hline
\end{tabular}}\par}

3. 이사비 \par
{\small\centering\renewcommand{\arraystretch}{1.25}
\resizebox{\ifdim\width>\linewidth\linewidth\else\width\fi}{!}{
\begin{tabular}{|c|c|c|}
\hline
\makecell{주택연면적} & \makecell{이사비} & \makecell{비고} \\
\hline
\makecell{66m\textsuperscript{2}～99m\textsuperscript{2}미만} & \makecell{1,540,000원} & \makecell{노임, 차량운임, 포장비 등 포함} \\
\hline
\makecell{99m\textsuperscript{2}이상} & \makecell{1,790,000원} & \makecell{노임, 차량운임, 포장비 등 포함} \\
\hline
\end{tabular}}\par}

※ 연면적은 부속건축물을 합한 면적을 적용하기로 함 \par
< 자료 10 > 축산업 관련 보상평가 참고자료 \par
1. 축산업의 가축별 기준마리 수 \par
{\small\centering\renewcommand{\arraystretch}{1.25}
\resizebox{\ifdim\width>\linewidth\linewidth\else\width\fi}{!}{
\begin{tabular}{|c|c|c|c|}
\hline
\makecell{가축} & \makecell{기준마리 수} & \makecell{가축} & \makecell{기준마리 수} \\
\hline
\makecell{소} & \makecell{5마리} & \makecell{닭} & \makecell{200마리} \\
\hline
\makecell{사슴} & \makecell{15마리} & \makecell{토끼} & \makecell{150마리} \\
\hline
\makecell{염소ㆍ양} & \makecell{20마리} & \makecell{오리} & \makecell{150마리} \\
\hline
\makecell{꿀벌} & \makecell{20군} & \makecell{돼지} & \makecell{20마리} \\
\hline
\end{tabular}}\par}

※ 자료: 「공익사업을 위한 토지 등의 취득 및 보상에 관한 법률 시행규칙」[별표 3] \par
2. 꿀벌(양봉) \par
가. 축산이익은 최근 3년간의 평균소득을 토대로 연간 240,000원/군 임 \par
나. 수송비는 5,000원/군 이고, 이전손실은 벌통폐사에 따른 손실, 채밀능력 저하 등으로 인한 손실, 이전시 유실 및 치사에 따른 손실 등을 고려할 때 25,000원/군 임(단, 그 외 추가적인 손실은 없는 것으로 함) \par
3. 닭(산란계) \par
가. 축산이익은 최근 3년간의 평균소득을 토대로 연간 3,600원/수 임 \par
나. 수송비는 200원/수 이고, 이전손실은 산란율 저하로 인한 손실, 폐사율 증가로 인한 손실, 제반 사육경비 등을 고려할 때 1,300원/수 임(단, 그 외 추가적인 손실은 없는 것으로 함) \par
4. 기타 참고사항 \par
가. 도시근로자가구 월평균 가계지출비(3인 가구 기준)는 4,665,400원 임 \par
나. 휴업기간은 4개월로 함 \par
다. 영업이익감소액은 휴업기간에 해당하는 영업이익의 100분의 20으로 하되, 1천 만을 초과하지 못함 \par
< 자료 11 > 기타 사항 \par
1. 지역요인비교치 및 개별요인비교치는 소수점 넷째자리에서 반올림하여 셋째자리까지 표시함 \par
2. 그 밖의 요인 보정치는 표준지 기준 산정방식을 적용함 \par
3. 대상토지의 결정단가는 백원단위에서 반올림하여 천원단위까지 표시함 \par
4. 건축물의 감가수정은 정액법으로 하며, 경과연수는 연단위로 산정함 \par
5. 건축물의 단가는 백원단위에서 절사하여 천원단위까지 표시함 \par
\end{longquestion}
